En esta pregunta usted debe definir la propiedad de ser resistente a colisiones para funciones de hash criptográficas. En particular, al utilizar frases tales como ``no existe un algoritmo de tiempo polinomial tal que ...''  debe indicar claramente cuáles son los parámetros en función de los cuáles se mide el tiempo de tales algoritmos.


\paragraph{Solución.} Una función de hash se define como una familia de funciones $\{h_n\}_{n \in \mathbb{N}}$ tal que $h_n : \{0,1\}^* \to \{0,1\}^{p(n)}$, donde $p$ es un polinomio. El parámetro $n$ corresponde a un parámetro de seguridad que determina cuál es el largo de la salida. Decimos entonces que $\{h_n\}_{n \in \mathbb{N}}$ es resistente a colisiones si no existe un algoritmo (aleatorizado) que funciona en tiempo $O(n^c)$ para una constante $c$, y que dado $n \in \mathbb{N}$, construye un par de strings $u, v \in \{0,1\}^*$ tal que $u \neq v$ y $h_n(u) = h_n(v)$.
